\SourcePage{81}{86}

\section{Relation between spinors and 4-vectors}

A spinor $\zeta^{\alpha\dot\beta}$ with one dotted and one undotted index has $2\cdot2=4$ independent components, precisely the number possessed by a 4-vector. It is therefore clear that the spinor and the 4-vector furnish the same irreducible representation of the proper Lorentz group, so a definite correspondence must exist between their components.

To establish this correspondence, we begin with its three-dimensional analogue, bearing in mind that 3- and 4-spinors must behave alike under purely spatial rotations.

For a three-dimensional spinor $\psi^{\alpha\beta}$, the correspondence formulas (see III, \S~57) may be written as
\[
\begin{aligned}
a_x&=\frac12(\psi^{22}-\psi^{11})=\frac12(\psi^2_1+\psi^1_2),\\
a_y&=-\frac{i}{2}(\psi^{22}+\psi^{11})=\frac{i}{2}(\psi^1_2-\psi^2_1),\\
a_z&=\frac12(\psi^{12}+\psi^{21})=\frac12(\psi^1_1-\psi^2_2),
\end{aligned}
\]
where $a_x$, $a_y$, and $a_z$ are the components of a certain three-dimensional vector $\mathbf a$. In passing to four dimensions, the components $\psi^\alpha_{\beta}$ must be replaced by $\zeta^{\alpha\dot\beta}$, while $a_x$, $a_y$, and $a_z$ are to be understood as the contravariant components $a^1$, $a^2$, and $a^3$ of a 4-vector. The form of the fourth component, $a^0$, follows at once from the circumstance noted in \S~17: quantity (17.6) must transform as $a^0$. Thus $a^0\sim\zeta^{1\dot1}+\zeta^{2\dot2}$; the proportionality coefficient is chosen so that the scalar $\zeta_{\alpha\dot\beta}\zeta^{\alpha\dot\beta}$ equals the scalar $2a_\mu a^\mu=2a^2$.

\SourcePage{82}{87}
The resulting correspondence formulas are
\begin{equation}
\begin{aligned}
a^1&=\frac12(\zeta^{1\dot2}+\zeta^{2\dot1}),
&a^2&=\frac{i}{2}(\zeta^{1\dot2}-\zeta^{2\dot1}),\\
a^3&=\frac12(\zeta^{1\dot1}-\zeta^{2\dot2}),
&a^0&=\frac12(\zeta^{1\dot1}+\zeta^{2\dot2}).
\end{aligned}
\tag{18.1}
\end{equation}
The inverse formulas are
\begin{equation}
\begin{aligned}
\zeta^{1\dot1}=\zeta_{2\dot2}&=a^3+a^0,
&\zeta^{2\dot2}=\zeta_{1\dot1}&=a^0-a^3,\\
\zeta^{1\dot2}=-\zeta_{2\dot1}&=a^1-ia^2,
&\zeta^{2\dot1}=-\zeta_{1\dot2}&=a^1+ia^2.
\end{aligned}
\tag{18.2}
\end{equation}
They give
\begin{equation}
\zeta_{\alpha\dot\beta}\zeta^{\alpha\dot\beta}=2a^2.
\tag{18.3}
\end{equation}
We also note that
\begin{equation}
\zeta_{\alpha\dot\beta}\zeta^{\gamma\dot\beta}=\delta_\alpha^\gamma a^2.
\tag{18.4}
\end{equation}
The last equality follows because the second-rank spinor $\zeta_{\alpha\dot\beta}\zeta^{\gamma\dot\beta}$ is antisymmetric in $\alpha\gamma$ and is therefore proportional to the metric spinor.

The correspondence between $\zeta^{\alpha\dot\beta}$ and a 4-vector is a special case of a general rule: every symmetric spinor of rank $(k,k)$ is equivalent to a symmetric irreducible 4-tensor of rank $k$, where irreducibility here means that it vanishes upon simplification over any pair of indices.

The relation between a spinor and a 4-vector can be written compactly using the two-row Pauli matrices\footnote{To simplify the notation, operators (matrices) acting on spin variables will be denoted by letters without hats.}:
\begin{equation}
\sigma_x=\begin{pmatrix}0&1\\1&0\end{pmatrix},\qquad
\sigma_y=\begin{pmatrix}0&-i\\i&0\end{pmatrix},\qquad
\sigma_z=\begin{pmatrix}1&0\\0&-1\end{pmatrix}.
\tag{18.5}
\end{equation}
Let $\zeta$ symbolically denote the matrix of the quantities $\zeta^{\alpha\dot\beta}$ with upper indices, the first being undotted. Formulas (18.2) then take the form
\begin{equation}
\zeta=\mathbf a\boldsymbol\sigma+a^0
\tag{18.6}
\end{equation}
\SourcePage{83}{88}
(the second term is, of course, understood to be the product of $a^0$ and the unit matrix). The inverse formulas are
\begin{equation}
\mathbf a=\frac12\Tr(\zeta\boldsymbol\sigma),
\qquad
a^0=\frac12\Tr\zeta.
\tag{18.7}
\end{equation}

Formulas (18.6) and (18.7) allow the transformation laws of a 4-vector and a spinor to be related. The spinor transformation law can thereby be expressed in terms of the rotation parameters of the four-dimensional coordinate system.

Write the transformation of $\xi^\alpha$ as
\begin{equation}
\xi^{\alpha'}=(B\xi)^\alpha,
\qquad
B=\begin{pmatrix}\alpha&\beta\\\gamma&\delta\end{pmatrix},
\tag{18.8}
\end{equation}
where $B$ is the two-row matrix formed from the coefficients of the binary transformation. The dotted spinor then transforms according to
\begin{equation}
\eta^{\dot\beta'}=(B^*\eta)^{\dot\beta}=(\eta B^+)^{\dot\beta},
\tag{18.9}
\end{equation}
and the transformation of the second-rank spinor $\zeta^{\alpha\dot\beta}\sim\xi^\alpha\eta^{\dot\beta}$ can be written symbolically as $\zeta'=B\zeta B^+$\footnote{For the covariant components,
\begin{equation}
\xi'_\alpha=(\widetilde B^{-1}\xi)_\alpha=(\xi B^{-1})_\alpha,
\qquad
\eta'_{\dot\alpha}=(\eta B^{*-1})_{\dot\alpha}
\tag{18.8a}
\end{equation}
(so that the product of two spinors $\xi_\alpha\Xi^\alpha$ remains invariant).}. For an infinitesimal transformation $B=1+\lambda$, where $\lambda$ is a small matrix, retention of terms through first order gives
\begin{equation}
\zeta'=\zeta+(\lambda\zeta+\zeta\lambda^+).
\tag{18.10}
\end{equation}

First consider a Lorentz transformation to a reference frame moving with infinitesimal velocity $\delta\mathbf V$, with no change in the directions of the spatial axes. The 4-vector $a^\mu=(a^0,\mathbf a)$ then transforms as
\begin{equation}
\mathbf a'=\mathbf a-a^0\delta\mathbf V,
\qquad
a^{0'}=a^0-\mathbf a\delta\mathbf V.
\tag{18.11}
\end{equation}
We now use (18.7). On the one hand, the transformation of $a^0$ may be written
\[
a^{0'}=a^0-\mathbf a\delta\mathbf V
=a^0-\frac12\Tr(\zeta\boldsymbol\sigma\delta\mathbf V),
\]
and, on the other hand,
\[
a^{0'}=\frac12\Tr\zeta'
=a^0+\frac12\Tr(\lambda\zeta+\zeta\lambda^+)
=a^0+\frac12\Tr\zeta(\lambda+\lambda^+).
\]

\SourcePage{84}{89}
These expressions must agree identically, that is, for arbitrary $\zeta$. It follows that
\[
\lambda+\lambda^+=-\boldsymbol\sigma\delta\mathbf V.
\]
Applying the same procedure to the transformation of $\mathbf a$ gives
\[
\boldsymbol\sigma\lambda+\lambda^+\boldsymbol\sigma=-\delta\mathbf V.
\]
Regarded as equations for $\lambda$, these relations have the solution
\[
\lambda=\lambda^+=-\frac12\boldsymbol\sigma\delta\mathbf V.
\]

Thus the infinitesimal Lorentz transformation of $\xi^\alpha$ is effected by the matrix
\begin{equation}
B=1-\frac12(\boldsymbol\sigma\mathbf n)\,\delta V,
\tag{18.12}
\end{equation}
where $\mathbf n$ is a unit vector in the direction of the velocity $\delta\mathbf V$. The transformation for a finite velocity $\mathbf V$ is now readily obtained. Recall that, geometrically, a Lorentz transformation is a rotation of the four-dimensional coordinate system in the $t\mathbf n$ plane through an angle $\varphi$, related to $\mathbf V$ by $\tanh\varphi=V$\footnote{Recall that the metric is pseudo-Euclidean in planes containing the time axis.}. An infinitesimal transformation corresponds to the angle $\delta\varphi=\delta V$, and a finite rotation through $\varphi$ is obtained by repeating the rotation through $\delta\varphi$ a total of $\varphi/\delta\varphi$ times. Raising the operator (18.12) to the power $\varphi/\delta\varphi$ and then taking the limit $\delta\varphi\to0$, we find
\begin{equation}
B=e^{-\frac{\varphi}{2}\mathbf n\boldsymbol\sigma}.
\tag{18.13}
\end{equation}
The mathematical content of this operator becomes clear from the properties of the Pauli matrices: every even power of $\mathbf n\boldsymbol\sigma$ is 1, whereas every odd power is $\mathbf n\boldsymbol\sigma$. Since the expansion of $\cosh$ contains the even powers of its argument and that of $\sinh$ the odd powers, the final result is
\begin{equation}
B=\cosh\frac{\varphi}{2}
-\mathbf n\boldsymbol\sigma\sinh\frac{\varphi}{2},
\qquad
\tanh\varphi=V.
\tag{18.14}
\end{equation}
The Lorentz-transformation matrices $B$ are thus Hermitian: $B=B^+$.

We next consider an infinitesimal rotation of the spatial coordinate system. The three-dimensional vector $\mathbf a$ transforms as
\begin{equation}
\mathbf a'=\mathbf a-\delta\boldsymbol\theta\times\mathbf a,
\tag{18.15}
\end{equation}

\SourcePage{85}{90}
where $\delta\boldsymbol\theta$ is the vector of the infinitesimal rotation angle. The corresponding spinor transformation could be found in the same way. This is unnecessary, however, because 4-spinors and 3-spinors behave alike under spatial rotations, while for the latter the transformation is already known from the general relation between the spin operator and the generator of an infinitesimal rotation:
\begin{equation}
B=1+\frac{i}{2}\boldsymbol\sigma\delta\boldsymbol\theta.
\tag{18.16}
\end{equation}
Passage to a finite rotation through an angle $\theta$ proceeds as in the passage from (18.12) to (18.14):
\begin{equation}
B=\exp\left(\frac{i\theta}{2}\mathbf n\boldsymbol\sigma\right)
=\cos\frac{\theta}{2}+i\mathbf n\boldsymbol\sigma\sin\frac{\theta}{2},
\tag{18.17}
\end{equation}
where $\mathbf n$ is a unit vector along the rotation axis. This matrix is unitary ($B^+=B^{-1}$), as is required for a spatial rotation.
